\begin{abstract}
Large language model (LLM) agents are increasingly evaluated on \emph{retrospective} memory---retrieving facts from long interaction histories to answer questions. Reliable assistants, however, must also maintain \emph{prospective} memory: deferred intentions that should be executed when a future cue or clock condition becomes true, while the agent continues other mandatory activities, and that must be updated under reschedule, cancel, and override events without spurious false alarms. Recent work on PM-Bench~\cite{pmbench2026} shows that this ``when-to-act'' competence is far from solved: monitoring hidden environment channels is a severe bottleneck, cross-day and update-sensitive hits remain low, and no universal scaffold dominates the Set-F1 operating point. We argue that the missing abstraction is \emph{conditional memory}: structured records of the form $(\varphi,\alpha,\sigma)$ that couple a trigger condition $\varphi$, an executable action $\alpha$, and a lifecycle state $\sigma$, rather than free-form episodic snippets optimized for needle retrieval. We formalize this view and propose \textbf{ProMem}, a prospective-memory scaffold with (i)~dual stores separating an episodic retrospective bank from a Prospective Intention Store (PIS), (ii)~an explicit lifecycle manager over announce$\rightarrow$active$\rightarrow$fired/completed/canceled/modified transitions, (iii)~a proactive trigger monitor that decides when to query clocks, email, and other hidden channels, and (iv)~a due-set scorer that selects executable candidates while suppressing lures. On the PM-Bench Virtual Week protocol we expect ProMem to improve Set-F1 especially on channel-required, cross-day, and update-sensitive slices, and to achieve a better precision--recall trade-off than spammy auto-heartbeats. Code and prompts will be released upon publication.
\end{abstract}
